\documentclass[10pt,a4paper]{amsart}
\usepackage[margin=19mm]{geometry}
\usepackage{amsmath,amssymb,mathtools}
\usepackage[scr=rsfs,cal=euler]{mathalfa}
\usepackage{xfrac}
\usepackage[unicode,hidelinks,bookmarks=false]{hyperref}
\newcommand{\ie}{i.e.\ }
\newcommand{\cf}{cf.\ }
\newcommand{\resp}{respectively\ }
\newcommand{\Subsection}{\subsection}
\providecommand{\nobreakdash}{\nobreak\hbox{-}}
\setlength{\emergencystretch}{2em}
\multlinegap0pt
\usepackage{bm}
\usepackage{bbm}
\usepackage{dsfont}
\usepackage{xspace}
\usepackage{cancel}
\usepackage{mathtools}
\usepackage{amsthm}
\makeatletter
\@ifundefined{lemma}{\newtheorem{lemma}{Lemma}}{}
\makeatother
\DeclareMathOperator{\E}{\mathbb{E}}
\title[Outlier-Robust Linear System Identification Under Heavy-tail]{Outlier-Robust Linear System Identification Under Heavy-tailed Noise}
\author{Vinay Kanakeri, Aritra Mitra}
\date{Source: arXiv:2501.00421v2 (CC BY 4.0)}
\begin{document}
\maketitle

\noindent\emph{Source and licence.} Outlier-Robust Linear System Identification Under Heavy-tailed Noise. Version arXiv:2501.00421v2 (CC BY 4.0), distributed under the Creative Commons Attribution 4.0 International licence, as recorded in the arXiv licence metadata for this exact version and shown on the abstract page https://arxiv.org/abs/2501.00421v2. Full rights notice: licenses/arxiv-2501.00421-CC-BY-4.0.txt.\par\smallskip

\noindent\textit{Editorial scope.} Counted unit: the Lemma whose source label is \texttt{None} in the article (source0.tex lines 549--559, source label \texttt{None}), delivered together with the article's own complete original proof of it (source0.tex lines 561--607, one \texttt{proof} environment of 47 lines). No mathematics is written, repaired or re-derived: statement and proof are the article's own text line for line. Required context retained uncounted: eqn:sys\_model (lines 149--152); eqn:noise\_model (lines 154--157). Every definition, display and label of those lines is retained unchanged. Omitted from this package and not redistributed: 1--148, 153--153, 158--548, 560--560, 608--765. Nothing inside a retained excerpt is omitted or altered: every retained body is delivered exactly as the article's lines are written, so no omission receipt of any kind accompanies this package. Citations: the retained excerpts carry no citation command at all, so no bibliography entry of the article is transcribed and no reference is redistributed. Reference adaptation: none was needed and none was made; every reference inside the delivered unit resolves inside it, so no unresolved reference or citation is produced. Printed slips kept literal: no printed word, symbol, hypothesis or number of the counted statement or of its proof was altered. Layout and class: the article's own class is \texttt{article} with packages including amsmath, amsfonts, amssymb, color, amsthm, nccmath, epsfig, psfrag, fontenc, algorithm, algpseudocode, xcolor, array, epstopdf; the wrapper is the lane's pinned \texttt{amsart} editorial wrapper at 10pt on A4, which restates the theorem-like environments the retained text needs and adds the article's own macro definitions that the retained text needs (\texttt{\textbackslash E}), with extra wrapper packages (bm, bbm, dsfont, xspace, cancel, mathtools). The delivered numbering is the unit's own. Assets: no retained span contains a figure environment, a graphics-inclusion command, a PDF-page inclusion, a table or a TikZ picture; no image file is copied into this package, the package declares no asset dependency and no image file is redistributed with it. Licence: version arXiv:2501.00421v2 of this article is distributed under the Creative Commons Attribution 4.0 International licence, as recorded in the arXiv licence metadata for this exact version and shown on the abstract page https://arxiv.org/abs/2501.00421v2. A full rights notice is delivered at licenses/arxiv-2501.00421-CC-BY-4.0.txt. Characters: every retained excerpt is pure ASCII, so the delivered engine, XeLaTeX with its default Latin Modern text font, reports no missing character.
\begin{equation}
x_{t+1}= A x_t + w_t,
\label{eqn:sys_model}
\end{equation}

\begin{equation}
\mathbb{E}[w_t w^{\top}_t] = \sigma^2 I_d, \hspace{3mm} \mathbb{E}[(w_t(i))^4] = \tilde{\sigma}^4, \forall i \in [d], \forall t \geq 0. 
\label{eqn:noise_model}
\end{equation}

\begin{lemma} Define $n_t \triangleq w_{T-(t+1)}.$ Given the system in \eqref{eqn:sys_model} and the noise assumptions in \eqref{eqn:noise_model}, we have
    \begin{equation}
    \resizebox{0.8\hsize}{!}{$
    \begin{aligned}
\E\left[(x_Tx_T^{\top})^2\right] &= \sum_{t=0}^{T-1} \E\left[(A^tn_tn_t^{\top}(A^t)^{\top})^2\right] 
         + 2\sum_{s\neq t=0}^{T-1} \E\left[ A^tn_tn_t^{\top}(A^t)^{\top} A^sn_sn_s^{\top}(A^s)^{\top}\right]\\ 
         &+ \sum_{s\neq t=0}^{T-1} \E\left[ A^tn_tn_s^{\top}(A^s)^{\top} A^sn_sn_t^{\top}(A^t)^{\top}\right]. 
         \nonumber 
\end{aligned}$}
    \end{equation}
\end{lemma}

\begin{proof}
    We have $x_T = \sum_{t = 0}^{T-1} A^tn_t$. In what follows, we simplify the summation notation by dropping the ranges, which vary from $0$ to $T-1$, and preserving only the relation between the indices for clarity. $\E\left[(x_Tx_T^{\top})^2\right]$ can be expressed as
    \begin{align*}
        \E\left[(x_Tx_T^{\top})^2\right] &= \E\left[ \left( \sum_{k}A^kn_k \sum_{l}n_l^{\top}(A^l)^{\top} \right) \left( \sum_{s}A^sn_s \sum_{t}n_t^{\top}(A^t)^{\top} \right)\right] \\
        &= \E \left( \underbrace{\sum_k A^kn_kn_k^{\top}(A^k)^{\top}}_{T_1} + \underbrace{\sum_{k\neq l} A^kn_kn_l^{\top}(A^l)^{\top}}_{T_2} \right) \\
        & \times \left( \underbrace{\sum_s A^sn_sn_s^{\top}(A^s)^{\top}}_{T_3} + \underbrace{\sum_{s\neq t} A^sn_sn_t^{\top}(A^t)^{\top}}_{T_4} \right).
    \end{align*}
    Next, we expand the above expression, and analyze the result of taking the expectation on each of the four terms that arise in the above product, starting with the term $T_1 \times T_3.$ We have 
    \begin{align*}
        \E[T_1 \times T_3] = \sum_{k} \E\left[(A^kn_kn_k^{\top}(A^k)^{\top})^2\right] + \sum_{k\neq s} \E\left[ A^kn_kn_k^{\top}(A^k)^{\top} A^sn_sn_s^{\top}(A^s)^{\top}\right]. 
    \end{align*}
    In the following, as a result of the noise process being i.i.d. with zero mean, we show that the term $T_1 \times T_4$ does not survive the expectation.
    \begin{align*}
        \E[T_1 \times T_4] & = \sum_{s \neq t} \sum_k \E\left[ A^kn_kn_k^{\top}(A^k)^{\top} A^sn_sn_t^{\top}(A^t)^{\top}\right] \\
        & \overset{(a)}{=}  \sum_k \sum_{s \neq t = k} \E\left[ A^kn_kn_k^{\top}(A^k)^{\top} A^sn_sn_k^{\top}(A^k)^{\top}\right] \\
        & + \sum_k \sum_{t \neq s = k} \E\left[ A^kn_kn_k^{\top}(A^k)^{\top} A^kn_kn_t^{\top}(A^t)^{\top}\right] \\
        & + \sum_k \sum_{t \neq s \neq k} \E\left[ A^kn_kn_k^{\top}(A^k)^{\top} A^sn_sn_t^{\top}(A^t)^{\top}\right] = 0.
    \end{align*}
    Note that all the three expectations in (a) evaluate to a $0$ matrix. For example, in the first term of (a), the embedded scalar term $n_k^{\top}(A^k)^{\top} A^sn_s$ can be moved to the right end, yielding: 
    $$ \E\left[ A^kn_kn_k^{\top}(A^k)^{\top} A^sn_sn_k^{\top}(A^k)^{\top}\right] = \E\left[ A^kn_k n_k^{\top}(A^k)^{\top} n_k^{\top}(A^k)^{\top}\right]  \E\left[ A^sn_s \right] =0,$$
where we used the fact that $n_k$ and $n_s$ are independent, and $\E[n_s]=0.$ A similar argument can be used to conclude that the second and third terms in (a) also evaluate to $0$. 

    Similarly, since $\E[T_1 \times T_4] = \E[T_2 \times T_3]$, we conclude that $\E[T_1 \times T_4] = 0$. Next, we analyze the term $\E[T_2 \times T_4]$ in the following:
    \begin{align*}
        \E[T_2 \times T_4] & = \sum_{k \neq l} \sum_{s \neq t} \E\left[A^kn_kn_l^{\top}(A^l)^{\top} A^sn_sn_t^{\top}(A^t)^{\top}\right] \\
        & \overset{(a)}{=} \sum_{k\neq l} \E\left[ A^kn_kn_k^{\top}(A^k)^{\top} A^ln_ln_l^{\top}(A^l)^{\top}\right] \\
        & + \sum_{k\neq l} \E\left[ A^kn_kn_l^{\top}(A^l)^{\top} A^ln_ln_k^{\top}(A^k)^{\top}\right] \\
        & \overset{(b)}{+} \sum_{l \neq k = t \neq s \neq l} \E\left[ A^kn_kn_l^{\top}(A^l)^{\top} A^sn_sn_k^{\top}(A^k)^{\top}\right] \\
        & + \sum_{l \neq k = s \neq t \neq l} \E\left[ A^kn_kn_l^{\top}(A^l)^{\top} A^kn_kn_t^{\top}(A^t)^{\top}\right] \\
        & + \sum_{k \neq l = t \neq s \neq k} \E\left[ A^kn_kn_l^{\top}(A^l)^{\top} A^sn_sn_l^{\top}(A^l)^{\top}\right] \\
        & + \sum_{k \neq l = s \neq t \neq k} \E\left[ A^kn_kn_l^{\top}(A^l)^{\top} A^ln_ln_t^{\top}(A^t)^{\top}\right] \\
        & + \sum_{k \neq l \neq s \neq t \neq k, s \neq k, t \neq l} \E\left[ A^kn_kn_l^{\top}(A^l)^{\top} A^sn_sn_t^{\top}(A^t)^{\top}\right].
    \end{align*}
    In the above display, it is not hard to verify that except the first and second terms of (a), all the remaining terms starting from (b) evaluate to zero. To see this, notice that the embedded scalar terms of the form $n^{\top}A^{\top} An$ can be flipped or moved within the expectation. Then, using the i.i.d assumption and independence, one can verify the aforementioned claim. For completeness, we show how this can be done for the first and second terms starting from (b); the rest follow similarly. For the first term, we have
$$ \E\left[ A^kn_k n_l^{\top}(A^l)^{\top} A^sn_sn_k^{\top}(A^k)^{\top}\right] =  \E\left[ A^kn_k n_k^{\top} (A^k)^{\top}\right] \E\left[ n_l^{\top}(A^l)^{\top} \right] \E\left[ A^sn_s  \right] =0,$$
where we used independence of $n_k, n_l,$ and $n_s$, and $\E\left[ n_s \right]=0.$ Similarly, for the second term,  
\begin{align*}
    \E\left[ A^kn_kn_l^{\top}(A^l)^{\top} A^kn_kn_t^{\top}(A^t)^{\top}\right] &= \E\left[ A^kn_k n_k^{\top}(A^k)^{\top} A^l n_l n_t^{\top}(A^t)^{\top}\right] \\
    &= \E\left[ A^kn_k n_k^{\top} (A^k)^{\top}\right] \E\left[ A^l n_l  \right] \E\left[ n_t^{\top}(A^t)^{\top} \right],
\end{align*}
which evaluates to $0$, since $\E \left [ n_l \right] =0.$ Proceeding similarly, we conclude that the terms starting from (b) do not survive the expectation. Therefore, we have
    \begin{align*}
        \E[T_2 \times T_4] & = \sum_{k=s \neq l=t} \E\left[ A^kn_kn_k^{\top}(A^k)^{\top} A^ln_ln_l^{\top}(A^l)^{\top}\right] \\
        & + \sum_{k=t \neq l=s} \E\left[ A^kn_kn_l^{\top}(A^l)^{\top} A^ln_ln_k^{\top}(A^k)^{\top}\right].
    \end{align*}
    Finally, compiling the results of $\E[T_1 \times T_3]$ and $\E[T_2 \times T_4]$, we obtain the desired claim in the lemma.
\end{proof}

\end{document}
